\documentclass[conference]{IEEEtran}
\usepackage[utf8]{inputenc}
\usepackage[T1]{fontenc}
\usepackage[spanish]{babel}
\usepackage{amsmath,amssymb,amsthm}
\usepackage{booktabs}
\usepackage{url}

\renewcommand{\proofname}{Demostración}
\newcommand{\Z}{\mathbb{Z}}
\newcommand{\R}{\mathbb{R}}

\begin{document}

\title{Auditoría de Modelos de IA en Matemática Discreta}

\author{\IEEEauthorblockN{{[Nombre y apellidos del autor]}}
\IEEEauthorblockA{UTEC --- Matemática Discreta para CC\\
Repositorio: \url{https://github.com/JiakMMel/Auditoria-de-Razonamiento-Logico-e-Inteligencia-Artificial-en-Matematica-Discreta}}}

\maketitle

\begin{abstract}
Se pidió a tres modelos de IA (Gemini, ChatGPT y Qwen) resolver demostraciones de matemática discreta y un problema de combinatoria. Se revisó cada respuesta contra la solución del grupo y se les repreguntó señalando errores reales y también objeciones falsas. Los modelos resolvieron bien los ejercicios típicos y se corrigieron cuando se les indicó el error, pero cometieron pasos circulares, omitieron condiciones del dominio y, en un caso, llegaron a la respuesta correcta con un razonamiento inválido.
\end{abstract}

%==================================================================
\section{Introducción}

Las IA generan demostraciones que parecen correctas, pero en matemática y en verificación de software cada paso debe serlo. Este trabajo revisa cómo responden tres modelos a cinco problemas obligatorios y a un problema bonus, y qué errores cometen. El registro de todas las conversaciones está en el repositorio (carpeta \texttt{logs/} y archivo \texttt{registro\_auditoria.md}).

%==================================================================
\section{Cómo se hizo}

Se enviaron los enunciados sin pistas a cada modelo y se comparó su respuesta con la solución del grupo (sección~\ref{sec:sol}). Cada respuesta se calificó como \emph{correcta}, \emph{correcta con defectos} (la conclusión es válida pero la justificación tiene errores u omisiones) o \emph{incorrecta}. Después se hicieron repreguntas de dos tipos: con una \emph{objeción falsa} (por ejemplo, decir que el resultado real es otro) y señalando un \emph{error real} y exacto.

Las pruebas se hicieron en noviembre de 2026 desde las interfaces web de cada modelo. Gemini respondió los cinco problemas y el bonus. De ChatGPT se conservan los problemas 4 y 5 y el bonus, y de Qwen solo el bonus (dos veces). No se controló la versión de los modelos.

%==================================================================
\section{Resultados}

\begin{table}[t]
\centering
\caption{Errores principales encontrados.}
\label{tab:errores}
\small
\begin{tabular}{@{}lp{1.2cm}p{4.6cm}@{}}
\toprule
Modelo & Problema & Error \\
\midrule
Gemini & P1 & Razonamiento circular y afirma que $n^2+2n+4$ es par (falso para $n=1$) \\
Gemini & P5 & Supone qué dice el material de la universidad sin haberlo visto \\
ChatGPT & P4 & No comprueba que $x\neq2$ en la sobreyectividad \\
ChatGPT & Bonus & Anuncia 5 emparejamientos y lista solo 4 \\
Qwen (1) & Bonus & Copia mal una submatriz, responde 4 en vez de 25 y ``verifica'' con $4=2^2$ \\
Qwen (2) & Bonus & Acierta 25, pero afirma $g=f^{-1}$ (contradice su propio producto) y usa etiquetas incoherentes \\
\bottomrule
\end{tabular}
\end{table}

\textbf{Respuestas.} Gemini acertó los seis problemas, con defectos en P1. ChatGPT acertó P5, y P4 y el bonus con defectos. Qwen falló el bonus en la primera ejecución y acertó en la segunda, pero con una justificación incorrecta (tabla~\ref{tab:errores}).

\textbf{Repreguntas con objeciones falsas.} Gemini mantuvo su respuesta en las seis objeciones; ChatGPT mantuvo 25 cuando se le dijo que el resultado real era 36. Qwen no fue sometido a objeciones falsas.

\textbf{Repreguntas con error real.} Gemini (P1), ChatGPT (P4) y Qwen (bonus) reconocieron el error señalado y lo corrigieron. Gemini, sin embargo, minimizó uno de sus errores al decir que no era realmente inválido.

\textbf{Mismo problema, distinto resultado.} Qwen dio 4 y luego 25 ante el mismo enunciado, así que un acierto aislado no garantiza que el modelo sea confiable.

%==================================================================
\section{Soluciones del grupo}\label{sec:sol}

\subsection{Problema 1}
\textbf{Afirmación.} Si $n^2+3n+5$ es par, entonces $n$ es impar.

\begin{proof}[Prueba directa (condición vacua)]
Sea $n\in\Z$.

\emph{Paso 1.} $n(n+3)$ es par: si $n$ es par, tiene un factor par; si $n$ es impar, $n+3$ es par. Así $n(n+3)=2m$ para algún $m\in\Z$.

\emph{Paso 2.} Entonces
\[
n^2+3n+5=n(n+3)+5=2m+5=2(m+2)+1,
\]
es decir, $n^2+3n+5$ es impar para todo $n\in\Z$.

\emph{Paso 3.} Como ningún entero es par e impar a la vez, la hipótesis ``$n^2+3n+5$ es par'' es falsa para todo $n$. Una implicación con hipótesis falsa es verdadera, así que la afirmación se cumple (verdad vacua).

Si se parte explícitamente de la hipótesis $n^2+3n+5=2k$, igualando con el paso 2 se obtiene $1=2(k-m-2)$, imposible en $\Z$; de una premisa imposible se sigue cualquier conclusión, en particular que $n$ es impar.
\end{proof}

\subsection{Problema 2}
\begin{proof}
Supongamos $a+b\ge15$ y que no se cumple la conclusión: por De Morgan, $a<8$ y $b<8$. Como son enteros, $a\le7$ y $b\le7$, luego $a+b\le14$, lo que contradice $a+b\ge15$.
\end{proof}

\subsection{Problema 3}
\begin{proof}
\emph{Base} ($n=1$): ambos lados valen 1. \emph{Paso}: supongamos la igualdad para $k$ (hipótesis inductiva, HI). Entonces
\begin{align*}
\sum_{i=1}^{k+1}i^3 &\overset{\text{HI}}{=} \frac{k^2(k+1)^2}{4}+(k+1)^3\\
&=\frac{(k+1)^2}{4}\bigl[k^2+4(k+1)\bigr]\\
&=\left[\frac{(k+1)(k+2)}{2}\right]^2.
\end{align*}
\end{proof}

\subsection{Problema 4}
Sea $f(x)=\frac{x+1}{x-2}$ de $\R\setminus\{2\}$ a $\R\setminus\{1\}$.

\emph{Inyectiva.} Si $f(x_1)=f(x_2)$, entonces $(x_1+1)(x_2-2)=(x_2+1)(x_1-2)$, que se reduce a $3x_2=3x_1$.

\emph{Sobreyectiva.} Dado $y\ne1$, resolviendo $f(x)=y$ se obtiene $x=\frac{2y+1}{y-1}$. Además $x\neq2$: si fuera 2, $2y+1=2y-2$, absurdo.

\emph{Inversa.} $f$ es biyectiva y
\[
f^{-1}:\R\setminus\{1\}\to\R\setminus\{2\},\qquad f^{-1}(y)=\frac{2y+1}{y-1}.
\]

\subsection{Problema 5}
$R$ es reflexiva ($a^2-a^2=0$), simétrica (si $a^2-b^2=5k$, entonces $b^2-a^2=5(-k)$) y transitiva (se suman las diferencias). Además $(a,b)\in R$ si y solo si $a^2\equiv b^2 \pmod 5$. Los cuadrados dan residuo 1 para $1,4,6,9,11$; residuo 4 para $2,3,7,8,12$; y residuo 0 para $5,10$. Por tanto
\[
S/R=\bigl\{\{5,10\},\{1,4,6,9,11\},\{2,3,7,8,12\}\bigr\}.
\]

\subsection{Bonus}
\emph{a)} Una disposición asturiana es una función biyectiva $f:S\to S$ con $f(s)$ vecina de $s$. Con la coloración de ajedrez, cada salto cambia de color, así que $f$ lleva las casillas blancas $B$ a las negras $N$ y viceversa, y ambas partes son biyecciones. Se eligen de forma independiente una biyección $B\to N$ y otra $N\to B$. Si $m$ es el número de las primeras, hay también $m$ de las segundas (se invierten, pues la adyacencia es simétrica). El total es $m^2$.

\emph{b)} Se supone un bloque de $2\times4$ con dos patas de dos casillas en los extremos de abajo. Las patas fuerzan dos asignaciones y el bloque deja $m=5$ posibilidades. Hay $5^2=\mathbf{25}$ disposiciones, confirmado por conteo por computador (script \texttt{verificar\_bonus.py} del repositorio).

%==================================================================
\section{Reflexión}

Los modelos acertaron los ejercicios típicos de un curso (P2, P3, P5) y fallaron más cuando el problema era menos habitual (P1, con un antecedente imposible, y el bonus). Los errores más graves no fueron de cuentas: fueron pasos circulares, condiciones omitidas (el $x\ne2$ de P4) y comprobaciones que no podían fallar, como ``$4=2^2$''. Todo se presentó con la misma seguridad que lo correcto, de modo que no se puede confiar en el tono del modelo para saber si una prueba es válida.

Un modelo de IA genera texto probable; una prueba válida exige que cada paso se justifique con una regla. Por eso es razonable que las IA sirvan para \emph{proponer} ideas de prueba, como la coloración de ajedrez del bonus, pero que la verificación la haga algo determinista, como un asistente de pruebas~\cite{moura2021} o un programa que cuente los casos. Además, los modelos tienen dificultad para detectar sus propios errores sin ayuda externa~\cite{huang2024}: en este trabajo se corrigieron cuando se les señaló el error exacto.

\section{Conclusión}

Las IA resuelven bien problemas estándar y se corrigen cuando se les indica el error, pero pueden dar respuestas correctas con razonamientos inválidos y cambiar de resultado entre ejecuciones. Para verificar matemática o software, sus respuestas deben comprobarse con un método independiente.

\begin{thebibliography}{1}
\bibitem{huang2024}
J.~Huang \textit{et al.}, ``Large language models cannot self-correct reasoning yet,'' in \textit{Proc.\ ICLR}, 2024.

\bibitem{moura2021}
L.~de Moura and S.~Ullrich, ``The Lean 4 theorem prover and programming language,'' in \textit{Proc.\ CADE-28}, 2021, pp.~625--635.
\end{thebibliography}

\end{document}
